% !TeX root = 20261020.tex
% !LW recipe=lablog XeLaTeX
\input{_support/style.tex}

\begin{document}

\reporttitle{2026-10-20}{연구 질문·현재 근거·후속 개선 방향}

\section*{목적과 이 보고서의 범위}

기존 연구 기록을 다시 읽기 위한 현재 연구 질문과 근거의 범위를 정리했다. 문서 개편일이며 새 모델 학습이나 실험을 실행한 날짜가 아니다.

\recordnote{수합·문서 개편: 2026-10-20. 실행 날짜는 각 절과 원시 실행 목록을 따른다. 아래의 과거 계획·진행 상태는 당시 기록이며, 현재 채택 상태와 최종 결과는 후속 실험을 함께 확인한다.}

\section{IDS 배경과 연구 질문}
IDS는 네트워크 연결의 특징을 보고 정상 통신과 공격, 또는 공격 유형을 예측하는 침입 탐지 시스템이다.
실제 운용에서는 기존에 관찰하지 않은 공격 사례가 라벨과 함께 추가될 수 있다.
따라서 고정된 학습·평가 집합의 분류 점수만으로 운용상의 유용성을 판단하기는 어렵다.
이번 연구는 소수 라벨 공격의 탐지 품질을 개선하는 것이 중심 목표이며,
공격 클래스를 순서대로 공개하는 실험은 그 목표를 확인하는 평가 시나리오다.

최근 순차 도입 실험은 기존 IDS 데이터의 실제 표본과 라벨을 사용해 공개 순서와 확보량을 구성했다.
평가 데이터 자체를 합성한 실험이 아니다. TabPFN 사전학습에 사용된 합성 과제와도 구별한다.
일부 이전 연구인 EXP65의 합성 사례 생성은 별도 실험이며, 최근 평가 데이터 전체가 합성이라는 뜻은 아니다.
순차 도입은 raw timestamp 트래픽을 그대로 재현한 실험이나 라벨 없는 미지 공격 탐지,
서로 다른 네트워크 간 일반화의 입증으로 확대 해석하지 않는다.

\section{사전학습, 미세조정, context 갱신}
TabPFN은 사전학습된 파라미터와 라벨이 있는 context를 함께 사용한다.
사전학습된 모델이 예측 방법을 갖고 있다는 것과 현재 IDS 데이터에서 각 클래스가 어떤 사례에 해당하는지
알고 있다는 것은 다르다. 예측은 개념적으로 $p_\theta(y\mid x,D_{\mathrm{context}})$이며,
가중치 $\theta$를 고정해도 $D_{\mathrm{context}}$는 필요하다.

현재 사용하는 v3 분류 구현에서는 context 라벨이 최종 클래스 확률 계산에 참여한다.
신규 공격 사례만 넣으면 기존 정상·공격 클래스를 같은 조건으로 분류할 근거가 사라진다.
새 사례를 추가하면서 기존 클래스의 사례도 유지하는 이유다.
미세조정은 손실을 역전파하여 가중치를 바꾸는 작업이고,
\path{fit_with_cache}는 현재 실험에서 전처리와 context 표현·캐시를 준비하는 작업이다.
이 두 작업을 모두 단순히 ``학습''이라고 부르면 비용과 설계가 혼동된다.

EXP77은 실제 초기 IDS 가중치 미세조정과 매 단계 미세조정을 비교했다.
반면 EXP83의 PFN은 초기 IDS 미세조정도 하지 않은 사전학습 가중치를 유지하고,
누적 또는 버퍼 context를 다시 구성한다. 이는 이미 캐시에 저장된 내용에 신규 행만 계산 비용 없이 덧붙이는 구현이 아니다.
따라서 가중치를 재학습하지 않는다는 사실만으로 갱신 비용이 0이거나 XGB 재학습보다 작다고 말할 수 없다.

\section{현재 사용할 평가 근거}
\begin{longtable}{p{.16\linewidth}p{.34\linewidth}p{.42\linewidth}}
\toprule 근거 & 평가 범위 & 사용 시 유의점 \\\midrule\endhead
EXP61--63 & 정제 데이터의 개발 holdout; 비교군, expert 역량, S/V 및 K 변화 & 단일 seed 탐색이며 방법별 정보 예산과 라벨 역할을 함께 표시한다. \\
EXP72--76 & 순차 도입 중 각 단계/서비스 버전의 예측을 합산 & 최종 모델 점수나 단계 F1의 산술평균이 아니다. \\
EXP77 & 초기 FT와 context 정책, 최종 모델 전체 test & 원본 query 순서 결과. FT 조건의 근거는 보존하되 최신 FT 없는 비교는 EXP81을 본다. \\
EXP79 & 저장된 대규모 모델의 최종 전체 test & 클래스가 묶인 query 순서의 진단 결과. DistPFN의 현재 비교 근거는 EXP81이다. \\
EXP81 & 소규모·10만 학습 사례의 최종 전체 test, 혼합 query & 모델 비교의 현재 근거. seed 43으로 행을 섞고 배치 8,192개를 사용했다. \\
EXP82 & ToN 제안 모델 복원·재구축·동일 context 대조 & 낮은 점수를 재현했고 Global 분할·expert 채택 제약을 분리해 진단했다. \\
EXP83 & 신규 공격당 10--400개, 누적·세 버퍼, 최종 전체 test & 작은 context의 품질·갱신 비용 비교. 비용 일부는 이전 실행 실측을 재사용했다. \\
\bottomrule
\end{longtable}

일반적인 Macro-F1은 클래스별 F1을 평균한 값이다. 여러 단계의 예측을 합산한 경우에는
그 예측으로 하나의 혼동행렬을 만든 후 F1을 계산하므로, 특정 클래스가 등장한 단계들의 F1 평균과 같지 않다.
이 문서의 최근 비교는 모든 클래스 도입이 끝난 마지막 모델로 전체 test를 예측한 값을 우선한다.

\section{현재 관찰과 다음 개선안}
신규 공격당 100개인 소규모 최종 평가에서 XGB/FT 없는 누적 PFN의 Macro-F1은
CIC2018 72.16/73.11\%, ToN 63.12/66.61\%다.
후속 갱신 합계는 각각 XGB/PFN 기준 2.61/1.26초, 5.62/1.54초였다.
반대로 누적 10만 사례에서는 PFN context 준비가 XGB 재학습보다 오래 걸렸다.
작은 context에서 확인한 비용 이점을 큰 context로 일반화할 수 없다.
초기 준비와 전체 추론은 이 갱신 비교에서 제외했으며, 운영 총비용의 우위는 별도 문제다.

제안 모델의 소규모 ToN 점수 52.89\%는 재구축해도 같았고, 동일한 Global 695개를 사용한 단일 PFN과 확률까지 일치했다.
전체 1,700개 context의 단일 PFN은 66.61\%였다.
따라서 지금은 Global에 남기는 라벨과 expert·검증 역할에 배정하는 라벨을 재설계할 근거가 있다.
적은 cal 표본에서 기존 최소 개입 근거를 충족하지 못한 점도 함께 개선해야 한다.
공통 context 활용, 역할별 교차적합, 충분한 교정 근거를 확보하는 검증 설계를 검토하되,
test 점수를 보고 호출 문턱을 임의로 낮추는 방식으로 해결하지 않는다.

EXP83에서는 신규 공격당 200·400개에서 누적 PFN이 세 버퍼보다 두 데이터 모두 높았다.
버퍼를 1,000개로 고정하고 신규 사례를 모두 넣으면 기존 정상·공격 사례가 과도하게 줄어드는 조건이 있다.
정상·기존 공격의 보존량과 신규 공격 배분을 함께 제어하고, 확보한 사례 중 대표 사례를 선택하는 방향이 다음 개선안이다.
모든 버퍼의 크기가 같아도 결과는 달랐으므로, 관찰된 차이를 크기만의 인과 효과라고 주장하지 않는다.
이 개선안들은 아직 실험으로 검증한 성과가 아니다.

\section{보존한 정정과 미채택 기록}
8월 13일의 일부 분포 변화 해석은 8월 17일에 정정되었다.
평가 표본 수를 클래스별로 제한한 과거 결과는 전체 test 결과와 구별한다.
EXP60은 실행했으나 현재 설계 근거에서 철회한 기록이다.
EXP64--67은 당시 별도 탐색 결과로, 일부 부분 표본 평가와 무효 실행을 포함한다.
EXP68은 사용자 의도와 다른 시나리오여서 채택하지 않는다.
과거 기록의 ``현재'', ``최고'', ``다음 단계'' 등 표현은 해당 기록 당시의 상태이며,
수합했다고 해서 그 주장이나 논문 근거 채택 상태가 현재로 갱신되는 것은 아니다.



\section*{출처와 실행 추적}

상세 수치의 원시 실행은 \path{inventory/runs.csv}, 실험별 연결은 \path{inventory/experiments.csv}에서 확인한다. 표·그림의 재현 입력은 \path{evidence/}와 \path{data/}에 둔다. 문서 이동·통합 내역은 \path{configs/experiments/migrations/20261020_documents/manifest.json}에 기록했다.

\end{document}
